% Part: methods
% Chapter: proofs
% Section: reading-proofs

\documentclass[../../../include/open-logic-section]{subfiles}

\begin{document}

\olfileid{mth}{prf}{rea}
\olsection{સાબિતીઓ વાંચવી}

પાઠ્યપુસ્તકો અને લેખોમાં મળતી સાબિતીઓમાં અત્યાર સુધીના
ઉદાહરણો જેવી બધી વિગતો ભાગ્યે જ લખાય છે. લેખકો ઘણી વાર
કિસ્સા અલગ પાડે છે, પરોક્ષ સાબિતી કરે છે, કે વ્યાખ્યા વાપરે છે
તે તરફ ધ્યાન દોરતા નથી. તેથી પાઠ્યપુસ્તકની સાબિતી વાંચતાં
તે સમજવા આવી વિગતો તમારે જાતે પૂરી કરવી પડે છે. આમ કરવાથી
સાબિતીમાં કરવાનાં વિવિધ પગલાંનો અભ્યાસ પણ થાય છે. એક
ઉદાહરણ જોઈએ.

\begin{prop}[શોષણ]
બધા ગણો $A$, $B$ માટે,
\[
A \cap (A \cup B) = A
\]
\end{prop}

\begin{proof}
જો $z \in A \cap (A \cup B)$, તો $z \in A$; તેથી $A \cap (A \cup B)
\subseteq A$. હવે $z \in A$ ધારો. ત્યારે $z \in A \cup B$ પણ છે;
તેથી $z \in A \cap (A \cup B)$ પણ છે.
\end{proof}

ઉપર શોષણના નિયમની સાબિતી ખૂબ સંક્ષિપ્ત છે. કઈ વ્યાખ્યાઓ
વપરાઈ છે તેનો ઉલ્લેખ નથી, સાબિત કરતા પહેલાં ``આ બતાવવાનું
છે'' એવું પણ લખાયું નથી. તેને ઉકેલીએ. સાબિત પ્રતિજ્ઞાપન
મનસ્વી ગણો $A$ અને $B$ અંગેનો સાર્વત્રિક દાવો છે; સાબિતીમાં
$A$ કે $B$ આવે ત્યારે તે મનસ્વી ગણો દર્શાવતા ચલો છે.
સાબિતી જે સામાન્ય દાવા સ્થાપિત કરે છે તે ગણોની સમાનતા
સાબિત કરવા જરૂરી છે: સમાનતાની ડાબી બાજુનું દરેક !!{element}
જમણી બાજુનું !!a{element} છે અને ઊલટું પણ.

\begin{quote}
``જો $z \in A \cap (A \cup B)$, તો $z \in A$; તેથી $A \cap (A \cup B)
  \subseteq A$.''
\end{quote}

આ સમાનતાની સાબિતીનો પહેલો ભાગ છે: ડાબી બાજુનું મનસ્વી~$z$
!!a{element} હોય, તો તે જમણી બાજુનું પણ !!a{element} છે,
એટલે $A \cap (A \cup B) \subseteq A$. $z \in A \cap (A \cup B)$
ધારો. બે ગણોના છેદગણમાં $z$ !!{element} હોય ત્યારે અને ત્યારે જ
જ્યારે તે બંને ગણોનું !!{element} હોય છે; તેથી $z \in A$
અને $z \in A \cup B$ મળે. ખાસ કરીને $z \in A$, જે બતાવવું
હતું. પહેલા ભાગ માટે આટલું પૂરતું હોવાથી બાકીની સાબિતી
બીજા ભાગની છે: $A \subseteq A \cap
(A \cup B)$ સાબિત કરવાની.

\begin{quote}
``હવે $z \in A$ ધારો. ત્યારે $z \in A \cup B$ પણ છે;
તેથી $z \in A \cap (A \cup B)$ પણ છે.''
\end{quote}

દરેક~$z$ માટે જો $z \in A$ તો $z \in A \cap (A \cup B)$
બતાવીએ છીએ, તેથી $z \in A$ ધારીને શરૂ કરીએ. $z
\in A \cap (A \cup B)$ બતાવવા $\cap$ની વ્યાખ્યા મુજબ
(i) $z \in A$ અને (ii) $z \in A \cup B$ બંને બતાવવું પડે.
અહીં (i) આપણી પૂર્વધારણા જ છે; વધુ કંઈ સાબિત કરવાનું નથી.
એટલા માટે સાબિતીમાં તેનો ફરી ઉલ્લેખ નથી. (ii) માટે યાદ કરો:
$z$ ગણોના યોગગણનું !!a{element} ત્યારે અને ત્યારે જ જ્યારે
તેમાંના ઓછામાં ઓછા એક ગણનું !!{element} હોય. $z \in A$
છે અને $A \cup B$ એ $A$ અને $B$નો યોગગણ છે, એટલે અહીં
આ શરત પૂરી થાય છે. તેથી $z \in A \cup B$. (i) $z \in A$
અને (ii) $z \in A \cup B$ બંને બતાવ્યાં; એટલે $\cap$ની
વ્યાખ્યા મુજબ $z \in A \cap (A \cup B)$. સાબિતી આ
વ્યાખ્યાઓનો ઉલ્લેખ કરતી નથી, કારણ કે વાચકને તે આવડી ગઈ છે
એમ માનવામાં આવે છે. ન આવડી હોય તો પાછા જઈને યાદ કરો.
પછી $z
\in A$ પરથી $z \in A \cup B$ કેમ મળે છે અને $z \in A$
તથા $z \in A \cup B$ પરથી $z \in A \cap (A \cup B)$
કેમ મળે છે તે પણ ઓળખવું પડે.

ઉપરની સાબિતીનું બધું સ્પષ્ટ કરતું બીજું સ્વરૂપ આ છે:
\begin{proof}{}
[$=$ની ગણો માટેની વ્યાખ્યા મુજબ $A \cap (A \cup B) = A$
બતાવવા (a) $A \cap (A \cup B) \subseteq A$ અને
(b) $A \subseteq A \cap (A \cup B)$ બતાવવું પડે.\sourcecorrection{OLMD-161}{મૂળના
(b)માં ભૂલથી (a) જેવી જ ઉપગણ-દિશા ફરી લખાઈ છે;
અહીં વિપરીત દિશા પૂરી છે.} (a): $\subseteq$ની વ્યાખ્યા
મુજબ જો $z \in A \cap (A \cup B)$ તો $z \in A$ બતાવવું પડે.]
જો $z \in A
\cap (A \cup B)$, તો $z \in A$ [$\cap$ની વ્યાખ્યા મુજબ
$z
  \in A \cap (A \cup B)$ ત્યારે અને ત્યારે જ જ્યારે
$z \in A$ અને $z \in A \cup B$], તેથી $A
\cap (A \cup B) \subseteq A$. [(b): $\subseteq$ની વ્યાખ્યા
મુજબ જો $z \in A$, તો $z \in A \cap (A \cup B)$ બતાવવું પડે.]
હવે [(1)] $z \in A$ ધારો. ત્યારે [(2)] $z \in A \cup B$ પણ છે
[(1) પરથી $z \in A$ અથવા $z \in B$, અને $\cup$ની વ્યાખ્યા
મુજબ તેનો અર્થ $z
  \in A \cup B$ થાય છે]; તેથી $z \in A \cap (A \cup B)$ પણ
છે [$\cap$ની વ્યાખ્યા મુજબ $z \in A$, એટલે (1), અને
$z
  \in A \cup B$, એટલે (2), બંને જોઈએ].\sourcecorrection{OLMD-162}{મૂળમાં
છેલ્લા યોગગણ-સભ્યતા સૂત્રમાં વધારાનું બંધ કૌંસ છે; અહીં
તે દૂર કર્યું છે.}
\end{proof}

\begin{prob}
$A \cup (A \cap B) = A$ની નીચેની સાબિતીને વિસ્તારો:
વપરાયેલા બધા અનુમાન-નમૂનાઓ જણાવો, દરેક પગલું પૂર્વધારણા
અથવા અગાઉ સ્થાપિત દાવામાંથી કેમ નીકળે છે તે કહો, અને
ક્યાં કઈ વ્યાખ્યાનો ઉપયોગ થાય છે તે બતાવો.
\begin{proof}
  જો $z \in A \cup (A \cap B)$, તો $z \in A$ અથવા
  $z \in A \cap B$. જો $z \in A \cap B$, તો $z \in A$.
  કોઈ પણ $z \in A$ એ $\in A \cup (A
  \cap B)$ પણ છે.
\end{proof}
\end{prob}

\end{document}
